
\documentclass[11pt]{article}
\usepackage{configurations}

\title{Unrolled Policy Iteration for Tiny Recursive Models:\\
Contraction, Bootstrapped Bellman Training, and Conservative Improvement}
\author{Anonymous}
\date{}

\begin{document}
\maketitle

\begin{abstract}
We cast Tiny Recursive Models (TRM)---which alternate inner latent updates and outer answer edits---as \emph{unrolled policy iteration}. The inner update $z\!\leftarrow\!f_\theta(z,y,x)$ plays the role of approximate policy evaluation, while the outer update $y\!\leftarrow\!g_\phi(y,z,x)$ performs policy improvement over \emph{edits} to the plan $y$. Our contributions are:
(i) a formal meta--MDP over plans; (ii) a stable, bootstrapped Bellman objective on an evaluator $V_\psi(z,x)$ using short rollouts and a target network; (iii) a stability tool: spectral/Lipschitz control of $f_\theta$ that guarantees inner-loop convergence as the number of unrolled evaluation steps grows; and (iv) monotone improvement guarantees for conservative edit policies (CPI/TRPO style) that explicitly account for evaluation error from finite unrolling and function approximation. The results yield actionable dials ($L_z,K,n,\delta$) with clear effects on stability and improvement.
\end{abstract}

\section{Motivation and setup}
\paragraph{TRM in brief.}
TRM performs $n$ inner updates of a latent $z\!\in\!\ZZ$ via a small network $f_\theta$, then updates an answer $y\!\in\!\YY$ via $g_\phi$, potentially repeating for $N_{\text{sup}}$ supervised steps. Empirically, recursion with deep supervision has shown strong generalization on symbolic tasks with small models~\citep{TRM2025}. We provide an RL lens and guarantees.

\paragraph{Meta--MDP over plans.}
Let $x\!\in\!\XX$ be an instance. A \emph{meta--state} is $s\!=\!(x,y)$ with current plan $y\!\in\!\YY$. An \emph{edit action} $a\!\in\!\mathcal{A}(y)$ transforms $y$ deterministically to $y'\!=\!\mathrm{edit}(y,a;x)$ (masking can restrict $\mathcal{A}(y)$). A task-specific \emph{checker} emits reward $r(s,a)\!\in\!\R$ (e.g., improvement in constraint satisfaction or pixel accuracy). With discount $\gamma\!\in\!(0,1)$, this defines a discounted MDP $(\SSS,\mathcal{A},P,r,\gamma)$ with $\SSS=\XX\times\YY$ and $P$ induced by the edit mechanism.\footnote{Our analysis does not require deterministic rewards; we allow stochastic rewards/transitions in the expectations.}

\paragraph{TRM as unrolled policy iteration.}
We maintain an internal \emph{evaluator state} $z\!\in\!\ZZ$ and a value head $V_\psi:\ZZ\times\XX\to\R$. The inner update
\begin{equation}
    z^{(t+1)} \;=\; f_\theta\!\bigl(z^{(t)},\,y,\,x\bigr),\qquad t=0,1,\dots,n-1
    \label{eq:inner}
\end{equation}
approximates policy evaluation for the current plan $y$, while the outer update
\begin{equation}
    y' \;=\; g_\phi\!\bigl(y,\,z^{(n)},\,x\bigr)
    \label{eq:outer}
\end{equation}
performs policy improvement (optionally via a stochastic edit policy $\pi_\phi(a\mid y,z,x)$). Training combines supervised losses on $y$ with RL losses on $(V_\psi,\pi_\phi)$.

\section{Operators and training objectives}

\subsection{Bellman operators on the meta--MDP}
Fix a (stochastic) edit policy $\pi(\cdot|s)$. The Bellman operator acts on $V:\SSS\to\R$ as
\begin{equation}
   (\mathcal{T}^\pi V)(s)\;=\;\E\!\left[r(s,a)\;+\;\gamma\,V(s') \;\middle|\; a\!\sim\!\pi(\cdot|s),\, s'\!\sim\!P(\cdot|s,a)\right].
\end{equation}
For $K\!\ge\!1$, define the $K$-step operator
\begin{equation}
 (\mathcal{T}^\pi_K V)(s)\;=\;\E\!\left[\sum_{k=0}^{K-1}\gamma^k r_k\;+\;\gamma^K V(s_K)\right],
 \label{eq:Kstep}
\end{equation}
where $(s_k,a_k,r_k)$ is the $k$-step rollout from $s$ under $\pi$.

\begin{lemma}[Contraction of $\mathcal{T}^\pi_K$]\label{lem:Kstep_contraction}
If $\gamma\in(0,1)$, then $\mathcal{T}^\pi_K$ is a $\gamma^K$-contraction in the sup norm:
$\|\mathcal{T}^\pi_K V-\mathcal{T}^\pi_K W\|_\infty \le \gamma^K \|V-W\|_\infty$.
Moreover, the unique fixed point is $V^\pi$.
\end{lemma}
\begin{proof}[Sketch]
Write $\mathcal{T}^\pi_K V=\sum_{k=0}^{K-1}\gamma^k R_k+\gamma^K P_\pi^K V$. Then
$\|\gamma^K P_\pi^K(V-W)\|_\infty\le \gamma^K\|V-W\|_\infty$. Banach's theorem yields uniqueness.
\end{proof}

\paragraph{Discounted occupancy.}
Let $d_\pi$ denote the \emph{normalized} discounted state distribution:
$d_\pi(s)=(1-\gamma)\sum_{t\ge 0}\gamma^t\,\mathbb{P}(s_t=s\mid \pi)$.
We use expectations $\E_{s\sim d_\pi}[\cdot]$ below.

\subsection{Bootstrapped Bellman training on $z$ via short rollouts}
Let $z^{(n)}(s)$ be the inner iterate after $n$ steps in~\eqref{eq:inner}, and let $U_n(s)\!:=\!V_\psi\!\bigl(z^{(n)}(s),x\bigr)$ approximate $V^\pi(s)$ for $s=(x,y)$. Using $K$-step returns with a slowly moving target network ($\bar\psi$) and a stop-gradient operator $\sg(\cdot)$ on the target, we minimize the MSE
\begin{equation}
\mathcal{L}_{\mathrm{BR}}(\psi)\;=\;\E_{s}\!\left[\;U_n(s)\;-\;\sg\!\Bigl(\textstyle\sum_{k=0}^{K-1}\gamma^k r_k+\gamma^K\, V_{\bar\psi}\!\bigl(z^{(n)}(s_K),x_K\bigr)\Bigr)\;\right]^2.
\label{eq:BR}
\end{equation}
The target network is updated by EMA: $\bar\psi \leftarrow \tau\bar\psi + (1-\tau)\psi$.

\begin{remark}[Target networks and ``double sampling'']
Eq.~\eqref{eq:BR} is a \emph{bootstrapped MSE} with a frozen (slow) target, not the residual-gradient objective of \citet{Baird1995Residual}. The stop-gradient $\sg$ treats the bootstrap as a label, which avoids the classical double-sampling pathology of squared Bellman residuals in the stochastic setting; short modelled rollouts further stabilize targets.
\end{remark}

\subsection{Conservative policy improvement over edits}
Let $\pi_\phi(a|s)$ be a stochastic edit policy and $A^\pi(s,a)=Q^\pi(s,a)-V^\pi(s)$ the advantage. We use the standard local surrogate
\begin{equation}
L_\pi(\pi') \;=\; \eta(\pi) + \frac{1}{1-\gamma}\,\E_{s\sim d_\pi,\,a\sim\pi'(\cdot|s)}\!\bigl[A^\pi(s,a)\bigr].
\label{eq:surrogate}
\end{equation}
A CPI/TRPO-style step either mixes $\pi$ and an improved $\pi_0$ via $\pi_{\text{new}}=(1-\alpha)\pi+\alpha\pi_0$ or constrains $\KL(\pi(\cdot|s)\,\|\,\pi'(\cdot|s))\le \delta$ on average. The canonical mixture bound is
\begin{equation}
\eta(\pi_{\text{new}})\;\ge\; L_\pi(\pi_{\text{new}})\;-\;\frac{2\epsilon\,\gamma}{(1-\gamma)^2}\,\alpha^2,
\qquad \epsilon \triangleq \max_{s}\E_{a\sim\pi_0(\cdot|s)}\bigl|A^\pi(s,a)\bigr|.
\label{eq:cpi_bound}
\end{equation}
Under a KL trust region, analogous bounds hold by relating KL to total variation (e.g., Pinsker).

\section{Stability tool: contraction of the inner evaluator}
We require that the inner map is a contraction in $z$.

\begin{assumption}[Uniform contraction in $z$ and completeness]\label{ass:Lip}
The metric space $(\ZZ,\norm{\cdot})$ is complete. There exists $L_z\in(0,1)$ such that for all $(x,y)\in\XX\times\YY$ and all $z,z'\in\ZZ$,
\begin{equation}
\norm{f_\theta(z,y,x)-f_\theta(z',y,x)} \;\le\; L_z\,\norm{z-z'}.
\end{equation}
\end{assumption}

\begin{remark}[How to enforce Assumption~\ref{ass:Lip}]
Assumption~\ref{ass:Lip} is the standard Banach contraction condition:
on a complete metric space $(\ZZ,\|\cdot\|)$, a uniformly $L_z$-Lipschitz
map in $z$ with $L_z<1$ has a unique fixed point for each $(x,y)$ and the
iteration $z^{(t+1)}=f_\theta(z^{(t)},y,x)$ converges linearly
(\cite{Denardo1967}). In neural implementations, using 1-Lipschitz
nonlinearities and bounding the product of spectral norms on the $z\!\to\!z$
path ensures $\Lip_z(f_\theta)<1$ (e.g., via spectral normalization;
\cite{Miyato2018}). We treat $L_z$ as a tunable stability dial, estimated
or upper-bounded during training.
\end{remark}

\begin{proposition}[Existence, uniqueness, and linear convergence of $z^*$]\label{prop:banach}
Under Assumption~\ref{ass:Lip}, for each $(x,y)$ the map $T_{x,y}(z)=f_\theta(z,y,x)$ has a unique fixed point $z^*(x,y)$ and the inner iteration~\eqref{eq:inner} converges linearly:
\begin{equation}
\norm{z^{(t)} - z^*(x,y)} \;\le\; L_z^{\,t}\,\norm{z^{(0)} - z^*(x,y)}.
\end{equation}
In particular, the bias due to finite unrolling is bounded by
$\norm{z^{(n)}-z^*}\le \frac{L_z^{\,n}}{1-L_z}\,\norm{z^{(1)}-z^{(0)}}$.
\end{proposition}

\paragraph{Spectral/Lipschitz control.}
If $f_\theta$ is a composition of 1-Lipschitz activations (e.g., ReLU) and linear layers $(W_\ell)_{\ell=1}^L$, then
$\Lip(f_\theta)\le\prod_{\ell=1}^L \|W_\ell\|_2$.
Enforcing spectral normalization $\|W_\ell\|_2\le s_\ell$ with $\prod_\ell s_\ell<1$ guarantees Assumption~\ref{ass:Lip}. The Lipschitz of the value head in $z$ is also a useful dial: if $\abs{V_\psi(z,x)-V_\psi(z',x)}\le L_V\norm{z-z'}$, then smaller $L_V$ directly tightens finite-unrolling error terms below.

\section{Error bounds and improvement guarantees}

\subsection{Bellman residual control and finite unrolling}
Define the \emph{used} value function $U_n(s)=V_\psi(z^{(n)}(s),x)$ and the \emph{fixed-point} value $U_*(s)=V_\psi(z^*(s),x)$.

\begin{proposition}[Two useful error bounds]\label{prop:error_bounds}
For any $K\ge 1$ and $\gamma\in(0,1)$:
\begin{enumerate}
\item[\textbf{(A)}] \textbf{Sharp bound at the used function.}
\begin{equation}
\|U_n - V^\pi\|_\infty \;\le\; \frac{1}{1-\gamma^K}\,\bigl\|U_n-\mathcal T_K^\pi U_n\bigr\|_\infty.
\label{eq:boundA}
\end{equation}
\item[\textbf{(B)}] \textbf{Separated sources (fixed-point residual + unrolling).}
If $V_\psi(\cdot,x)$ is $L_V$-Lipschitz in $z$,
\begin{equation}
\|U_n - V^\pi\|_\infty
\;\le\; \frac{1}{1-\gamma^K}\,\bigl\|U_*-\mathcal T_K^\pi U_*\bigr\|_\infty
\;+\; L_V\,\sup_{(x,y)}\|z^{(n)}(x,y)-z^*(x,y)\|.
\label{eq:boundB}
\end{equation}
\end{enumerate}
\end{proposition}
\begin{remark}
Eq.~\eqref{eq:boundA} ties training directly to the BR loss in Eq.~\eqref{eq:BR} (evaluate the residual at $U_n$). Eq.~\eqref{eq:boundB} attributes error to (i) the fixed-point residual (modeling/approximation) and (ii) finite unrolling; by Prop.~\ref{prop:banach}, the second term decays geometrically in $n$ when $L_z<1$.
\end{remark}

\subsection{Conservative improvement under evaluation error}
Let $\widehat{A}(s,a)$ be an estimator of $A^\pi(s,a)$ based on $V_\psi$ (e.g., $\lambda$-returns or $K$-step TD). Suppose we have a uniform error bound
$\abs{\widehat{A}(s,a)-A^\pi(s,a)}\le \varepsilon_A$.

\begin{assumption}[Centered estimator]\label{ass:centered}
$\E_{a\sim\pi(\cdot|s)}\bigl[\widehat{A}(s,a)\bigr]=0$ for all $s$ (e.g., by defining $\widehat A(s,a)=\widehat Q(s,a)-\sum_b \pi(b|s)\widehat Q(s,b)$).
\end{assumption}

\begin{theorem}[Monotone improvement up to evaluation error]\label{thm:monotone_with_error}
Let $\pi_{\text{new}}=(1-\alpha)\pi+\alpha\pi_0$. Under $\gamma\in(0,1)$ and Assumption~\ref{ass:centered},
\begin{equation}
\eta(\pi_{\text{new}})\;\ge\; \widehat{L}_\pi(\pi_{\text{new}})
\;-\;\frac{\alpha}{1-\gamma}\,\varepsilon_A
\;-\;\frac{2\epsilon\,\gamma}{(1-\gamma)^2}\,\alpha^2,
\qquad \epsilon \triangleq \max_{s}\E_{a\sim\pi_0(\cdot|s)}\abs{A^\pi(s,a)}.
\label{eq:main_theorem}
\end{equation}
In particular, if $\widehat{L}_\pi(\pi_{\text{new}})$ exceeds the two penalties on the right-hand side, the true return improves.
\end{theorem}
\begin{proof}[Sketch]
Start from~\eqref{eq:cpi_bound} and add/subtract $\widehat{A}$. The linear term introduces a bias bounded by $\frac{\alpha}{1-\gamma}\varepsilon_A$ under centering; the quadratic CPI penalty remains governed by the true $A^\pi$ via $\epsilon$.
\end{proof}

\begin{remark}[Without centering]
If Assumption~\ref{ass:centered} does not hold, a conservative variant is
\[
\eta(\pi_{\text{new}})\;\ge\; \widehat{L}_\pi(\pi_{\text{new}})
\;-\;\frac{1}{1-\gamma}\,\varepsilon_A
\;-\;\frac{2\epsilon\,\gamma}{(1-\gamma)^2}\,\alpha^2,
\]
i.e., the linear penalty no longer scales with $\alpha$.
\end{remark}

\begin{remark}[Relating $\varepsilon_A$ to training dials]
By Prop.~\ref{prop:error_bounds}, $\varepsilon_A$ can be bounded in terms of (i) the BR objective value (via Eq.~\eqref{eq:boundA} or \eqref{eq:boundB}) and (ii) the inner contraction dial $L_z$ and unrolling budget $n$. Spectral control (smaller $L_z$), larger $n$, and larger $K$ (tighter contraction factor $\gamma^K$) reduce these terms.
\end{remark}

\section{Algorithmic template}
\begin{algorithm}[H]
\caption{Unrolled Policy Iteration for TRM (single interaction step)}
\label{alg:upi}
\begin{algorithmic}[1]
\State \textbf{Input:} instance $x$, current plan $y$, latent $z$, discount $\gamma$, inner steps $n$, rollout horizon $K$, temperature/entropy coeff.\ $\beta$.
\State \textbf{Inner evaluation:} $z \leftarrow f_\theta^{\circ n}(z,y,x)$ \Comment{$n$ compositions}
\State \textbf{Policy improvement:} sample edit $a\sim\pi_\phi(\cdot\mid y,z,x)$; apply $y'\leftarrow \mathrm{edit}(y,a;x)$; receive $r$ and next state $s'=(x,y')$.
\State \textbf{$K$-step targets:} roll out $(s_k,a_k,r_k)_{k=0}^{K-1}$ under $\pi_\phi$; compute
\[
G^{(K)}=\sum_{k=0}^{K-1}\gamma^k r_k+\gamma^K V_{\bar\psi}\!\bigl(z^{(n)}(s_K),x_K\bigr).
\]
\State \textbf{Losses:}
\begin{align*}
\mathcal{L}_{\mathrm{BR}}(\psi) &= \bigl(V_\psi(z^{(n)}(s),x)-\sg\big(G^{(K)}\big)\bigr)^2,\\
\mathcal{L}_{\pi}(\phi) &= -\,\widehat{A}(s,a)\,\log \pi_\phi(a\mid y,z,x)\;-\;\beta\,\mathcal{H}\!\left[\pi_\phi(\cdot\mid y,z,x)\right],\\
\mathcal{R}_{\mathrm{Lip}}(\theta,\psi) &= \sum_{\ell}\bigl(\|W_\ell\|_2 - s_\ell\bigr)_+ \quad\text{(for both $f_\theta$ and, optionally, $V_\psi$)},\\
\mathcal{L}_{\mathrm{sup}}(\phi) &\;\text{if labels for } y \text{ are available.}
\end{align*}
\State \textbf{Update:} minimize $\alpha_{\mathrm{BR}}\mathcal{L}_{\mathrm{BR}}+\alpha_{\pi}\mathcal{L}_{\pi}+\alpha_{\mathrm{Lip}}\mathcal{R}_{\mathrm{Lip}}+\alpha_{\mathrm{sup}}\mathcal{L}_{\mathrm{sup}}$, with a conservative step on $\pi_\phi$ (e.g., mixture update or $\KL$ trust region / clipping).
\end{algorithmic}
\end{algorithm}

\section{Discussion: residuals vs.\ bootstrapped MSE}
Classical residual-gradient methods minimize a squared Bellman residual but suffer from the “double sampling’’ issue in stochastic settings~\citep{Baird1995Residual,SuttonBarto2018}.
Our training uses a \emph{bootstrapped MSE} with a target network (Eq.~\eqref{eq:BR}); the stop-gradient makes the bootstrap act like a label during backprop, and the EMA target ensures smooth label evolution. Short rollouts provide low-variance targets; Eq.~\eqref{eq:boundA}/\eqref{eq:boundB} shows why reducing the residual at the used function (or at $z^*$) is theoretically meaningful.

\section{Related work}
\textbf{Recursive reasoning with tiny networks.} TRM shows that tiny networks with recursive latent updates and deep supervision can match or beat much larger models on symbolic tasks~\citep{TRM2025}.\\
\textbf{Conservative improvement.} Our guarantee follows CPI~\citep{KakadeLangford2002} and TRPO~\citep{Schulman2015TRPO}.\\
\textbf{Residual learning and contraction.} Residual/Bellman-error methods~\citep{Baird1995Residual,SuttonBarto2018} and contraction arguments in dynamic programming~\citep{Denardo1967} underpin the evaluation side; spectral normalization supplies an implementable Lipschitz bound~\citep{Miyato2018}.

\section{Conclusion}
Unrolled policy iteration provides a principled lens on TRM: inner contraction (via spectral control) guarantees stable evaluation as $n$ grows; $K$-step, target-network training makes evaluation consistent; and conservative edit policies deliver monotone improvement up to explicit evaluation error. The theory turns into practical dials ($L_z,K,n,\delta$) with predictable effects on stability and improvement.

\bibliographystyle{plainnat}
\bibliography{trm_rl}
\end{document}